\section{Однобитный перенос правила на замороженной OLMo}
\label{sec:olmo}
В опыте Rule~110 после слияния в клетку возвращают исходное показание датчика $Z$. Автомат восстанавливает прежний бит $Y$ при следующем обновлении, используя оставшийся в решётке бит $K$. В двухшаговом опыте OLMo отправитель до предъявления новой строки сводит прошлый выбор $h$ и обратную связь $e$ к одному токену правила $m\in\{A,B\}$. Получатель с теми же замороженными весами начинает новый эпизод со свежим KV-кэшем и применяет правило к строке. Такая граница эпизодов позволяет измерить перенос уже выбранного правила. Вопрос о внутренней коррекции стратегии в непрерывном состоянии модели требует отдельного опыта.

\subsection{Контрфактический квадрат}
Пусть ранее модель выбрала правило $h=0$ (A, взять первый бит) или $h=1$ (B, взять последний). Наблюдатель сообщает $e=0$ при верном выборе и $e=1$ при ошибке. Для каждого выбора доступны обе ветви обратной связи. До предъявления любой новой строки исправленное правило равно $p=h\oplus e$. После предъявления $\mathbf{x}=LwR$ правильный ответ равен $f_0(\mathbf{x})=L$ или $f_1(\mathbf{x})=R$. На разных концах две истории $h$ при одном $e,\mathbf{x}$ требуют противоположных действий. На равных концах обе требуют одного бита. Предъявление всех четырёх сочетаний концов препятствует стратегии, выучившей только связь ответа с одной буквой обратной связи.

Одна середина задаёт блок из 16 условий $(L,R,h,e)$. В dev восемь середин, отобранных до измерения. Блок засчитывается только при верных 16/16 ответах. Порог допуска равен 6/8 блоков. Во всех обязательных числовых чтениях требуется $P(0)+P(1)\ge0.80$ по полному словарю. Строгий знак зазора $\ell_1-\ell_0$ определяет выбранный бит; нулевой зазор считается воздержанием. Строки с одним и тем же префиксом переиспользуют вычисление модели.

\subsection{Что было измерено}
В G0 кэш после обработки выбранного токена изолируется до всякого диагностического чтения. В G0R позиция выбора подменяется значениями из другой натуральной истории и затем восстанавливается. В исходном опыте Base все 256/256 сравнений полных логитов для двух направлений дали равенство исходных байтов и максимальную абсолютную разность ноль. Выбор противоположной стратегии этим контролем не проверяется.

\begin{center}\small
\begin{tabularx}{\linewidth}{@{}Xrrrl@{}}\toprule
Чекпойнт и интерфейс (dev) & Строки & Блоки & Пары при $e=1$ & Статус \\\midrule
Base, исходный raw & 76/128 & 0/8 & 0/16 & \texttt{FAIL\_DEV}\\
Instruct, выровненный raw & 79/128 & 0/8 & 1/16 & \texttt{FAIL\_FORMAT}\\
Instruct, chat, correct/wrong & 115/128 & 2/8 & 9/16 & \texttt{FAIL\_DEV}\\
Instruct, chat, keep/switch & 109/128 & 0/8 & 6/16 & \texttt{FAIL\_DEV}\\
Instruct, chat, два шага (A) & 124/128 & 6/8 & 16/16 & \texttt{PASS\_DEV}\\
\bottomrule\end{tabularx}
\end{center}
Пара в таблице требует двух верных ответов для противоположных историй при одних $e,\mathbf{x}$ и $L\ne R$. В двух первых строках вместе с весами изменились демонстрации. В raw Instruct 28/64 контрольных строк с явно заданным правилом нарушили порог вероятностной массы. При нативном chat-шаблоне с correct/wrong на отдельных копиях кэша правильно прочитаны 2/2 прежних выбора. При явном правиле получены 62/64 правильные строки и 6/8 блоков. \foreignlanguage{english}{Uzay~Macar} и соавт.~\cite{macar2026} при введении вектора в активации модели отдельно измеряли способность обнаружить вмешательство и способность назвать введённое понятие. Здесь поведенческий критерий требует применить полученную информацию к новой строке. Активационный вариант G3 сохраняет статус \texttt{NOT\_RUN}.

\subsection{Теорема об однобитном канале}
В двухшаговом варианте отправитель выдаёт уникальный максимум по полному словарю до появления $\mathbf{x}$. Допустимы только токены A и B. Получатель стартует с пустым собственным KV-кэшем. Он принимает лишь ID токена $m$, заново вычисляет его представление и затем получает $\mathbf{x}$. Идентификатор ветви, активации и логиты отправителя ему недоступны. Внешняя программа не выбирает код по меткам правила.

\begin{proposition}
Пусть $m$ --- единственный переменный вход получателя из предыдущего эпизода, его алфавит двухэлементен, а канал имеет вид $(h,e)\longrightarrow m\longrightarrow F(m,\mathbf{x})$. Если на одном общем входе с $L\ne R$ все четыре истории решены точно, то в порядке $00,01,10,11$ таблица отправителя равна $0110$ или $1001$.
\end{proposition}
\begin{proof}
Два разных $p$ при данном $\mathbf{x}$ требуют разных ответов. Общее значение $m$ для таких историй предъявило бы получателю одинаковые входы при разных целях. Значит, два непустых множества сообщений не пересекаются. Поскольку сообщений ровно два, каждому $p$ соответствует единственное $m$. Возможны только две перестановки имён токенов.
\end{proof}

На dev наблюдалась таблица $A,B,B,A$. При каждом $e$ правильно решены обе истории во всех 16/16 парах с разными концами. Это 64/64 правильных ответа при доставке и 32/64 в независимо вычисленном контроле с постоянным A. Четыре решения отправителя переиспользуются на новых строках получателя. Следовательно, информация о классе правила присутствует в сообщении до появления строки. Для установления того, как она вычислена внутри отправителя, требуются иные вмешательства. Многомерное состояние могло бы передать оба бита $(h,e)$ и предоставить вычисление правила получателю. Источники измерений и контроли перечислены в приложении~\ref{app:gates}.
